\chapter{2012年北京卷（文）}

\section{单选题}

\begin{problem}
已知集合 \(A = \{x \in \R \mid 3x + 2 > 0\}\)，\(B = \{x \in \R \mid (x + 1)(x - 3) > 0\}\)，则 \(A \cap B =\)（\quad）

\choices
    {\((-\infty, -1)\)}
    {\(\left(-1, -\frac{2}{3}\right)\)}
    {\(\left(-\frac{2}{3}, 3\right)\)}
    {\((3, +\infty)\)}
\end{problem}

\begin{answer}
D
\end{answer}

\begin{solution}
由 \(3x+2>0\) 得 \(x>-\frac{2}{3}\)；由 \((x+1)(x-3)>0\) 得 \(x<-1\) 或 \(x>3\). 因而 \(A\cap B=(3,+\infty)\)，故选 D.
\end{solution}

\begin{problem}
在复平面内，复数 \(\frac{10\i}{3 + \i}\) 对应的点的坐标为（\quad）

\choices
    {\((1,3)\)}
    {\((3,1)\)}
    {\((-1,3)\)}
    {\((3,-1)\)}
\end{problem}

\begin{answer}
A
\end{answer}

\begin{solution}
将分式的分子、分母同乘 \(3-\i\)，得
\(\dfrac{10\i}{3+\i}=\dfrac{10\i(3-\i)}{(3+\i)(3-\i)}=\dfrac{30\i-10\i^2}{10}=1+3\i\).
所以该复数在复平面内对应的点为 \((1,3)\)，故选 A.
\end{solution}

\begin{problem}
设不等式组
\examdisplaycases{}{
0 \leqslant x \leqslant 2,\\
0 \leqslant y \leqslant 2
}

表示的平面区域为 \(D\). 在区域 \(D\) 内随机取一个点，则此点到坐标原点的距离大于 \(2\) 的概率是（\quad）

\choices
    {\(\frac{\pi}{4}\)}
    {\(\frac{\pi - 2}{2}\)}
    {\(\frac{\pi}{6}\)}
    {\(\frac{4 - \pi}{4}\)}
\end{problem}

\begin{answer}
D
\end{answer}

\begin{solution}
区域 \(D\) 是边长为 \(2\) 的正方形，面积为 \(4\). 在该正方形内，到原点距离不大于 \(2\) 的部分是半径为 \(2\) 的四分之一圆，面积为 \(\frac{1}{4}\pi\cdot2^2=\pi\). 因而到原点距离大于 \(2\) 的部分面积为 \(4-\pi\)，其边界不影响几何概率，所以所求概率为 \(\dfrac{4-\pi}{4}\)，故选 D.
\end{solution}

\begin{problem}
执行如图所示的程序框图，输出的 \(S\) 值为（\quad）

\texfigure[max width=0.70\linewidth]{img_repaint/2012/beijing_liberal/q4_fig1.tex}

\choices
    {\(2\)}
    {\(4\)}
    {\(8\)}
    {\(16\)}
\end{problem}

\begin{answer}
C
\end{answer}

\begin{solution}
按照框图，初始时 \(k=0\)，\(S=1\). 每次先计算 \(S\leftarrow S\cdot2^k\)，再计算 \(k\leftarrow k+1\)，循环条件为 \(k<3\). 三次循环依次得到
\((S,k)=(1,1)\)、\((2,2)\)、\((8,3)\).
此时 \(k<3\) 不成立，输出 \(S=8\)，故选 C.
\end{solution}

\begin{problem}
函数 \(f(x) = x^{\frac{1}{2}} - \left(\frac{1}{2}\right)^x\) 的零点个数为（\quad）

\choices
    {\(0\)}
    {\(1\)}
    {\(2\)}
    {\(3\)}
\end{problem}

\begin{answer}
B
\end{answer}

\begin{solution}
函数定义域为 \([0,+\infty)\). 在此区间上，\(\sqrt{x}\) 严格递增，而 \(-\left(\frac{1}{2}\right)^x\) 也严格递增，因此 \(f(x)\) 严格递增，至多有一个零点. 又
\(f\left(\frac{1}{2}\right)=\sqrt{\frac{1}{2}}-\left(\frac{1}{2}\right)^{\frac{1}{2}}=0\)，所以零点恰有一个，故选 B.
\end{solution}

\begin{problem}
已知 \(\{a_n\}\) 为等比数列，下面结论中正确的是（\quad）

\choices
    {\(a_{1} + a_{3} \geqslant 2a_{2}\)}
    {\(a_{1}^{2} + a_{3}^{2} \geqslant 2a_{2}^{2}\)}
    {\(\text{若 }a_{1}=a_{3}\text{，则 }a_{1}=a_{2}\)}
    {\(\text{若 }a_{3}>a_{1}\text{，则 }a_{4}>a_{2}\)}
\end{problem}

\begin{answer}
B
\end{answer}

\begin{solution}
设等比数列的公比为 \(q\)，则 \(a_2=a_1q\)，\(a_3=a_1q^2\). 对选项 B，有
\(a_1^2+a_3^2-2a_2^2=a_1^2(1+q^4-2q^2)=a_1^2(q^2-1)^2\geqslant0\)，所以 B 正确.

选项 A 不恒成立，例如取 \(a_1=-1\)，\(q=2\)，则 \(a_1+a_3=-5<-4=2a_2\). 选项 C 中取 \(a_1=1\)，\(q=-1\)，有 \(a_1=a_3=1\)，但 \(a_2=-1\ne a_1\). 选项 D 中取 \(a_1=1\)，\(q=-2\)，有 \(a_3=4>a_1\)，但 \(a_4-a_2=a_1q(q^2-1)=-6<0\)，即 \(a_4<a_2\). 故选 B.
\end{solution}

\begin{problem}
某三棱锥的三视图如图所示，该三棱锥的表面积是（\quad）

\texfigure[max width=0.58\linewidth]{img_repaint/2012/beijing_liberal/q7_fig1.tex}

\choices
    {\(28 + 6\sqrt{5}\)}
    {\(30 + 6\sqrt{5}\)}
    {\(56 + 12\sqrt{5}\)}
    {\(60 + 12\sqrt{5}\)}
\end{problem}

\begin{answer}
B
\end{answer}

\begin{solution}
按三视图中各顶点的对应位置记底面顶点为 \(B\)，\(C\)，\(D\)，并记顶点 \(A\) 在底边 \(BC\) 上的正投影为 \(H\). 由俯视图，
\(\angle BCD=90^\circ\)，\(BC=2+3=5\)，\(CD=4\)，所以
\(S_{BCD}=\frac{1}{2}\cdot5\cdot4=10\)，且
\(BD=\sqrt{5^2+4^2}=\sqrt{41}\).

正视图给出 \(BH=2\)，\(HC=3\)，以及 \(AH=4\). 因此
\(AB=\sqrt{BH^2+AH^2}=\sqrt{2^2+4^2}=2\sqrt{5}\)，
\(AC=\sqrt{HC^2+AH^2}=\sqrt{3^2+4^2}=5\)，并且
\(S_{ABC}=\frac{1}{2}\cdot BC\cdot AH=\frac{1}{2}\cdot5\cdot4=10\).

由三视图中 \(A\)、\(C\) 的侧向投影重合，而 \(CD\) 沿侧向，得 \(AC\perp CD\). 所以
\(S_{ACD}=\frac{1}{2}\cdot AC\cdot CD=\frac{1}{2}\cdot5\cdot4=10\).

俯视图中 \(H\) 与 \(D\) 的水平投影两条直角边分别为 \(HC=3\) 和 \(CD=4\)，故 \(AD\) 的水平投影长为
\(\sqrt{3^2+4^2}=5\). 结合正视图中 \(A\) 与底面的高差为 \(4\)，得
\(AD=\sqrt{5^2+4^2}=\sqrt{41}=BD\). 因而三角形 \(ABD\) 为等腰三角形，取 \(AB\) 为底边，其高为
\(\sqrt{AD^2-\left(\frac{AB}{2}\right)^2}=\sqrt{41-5}=6\)，从而
\(S_{ABD}=\frac{1}{2}\cdot2\sqrt{5}\cdot6=6\sqrt{5}\).

所以三棱锥的表面积为
\(S=S_{BCD}+S_{ABC}+S_{ACD}+S_{ABD}=10+10+10+6\sqrt{5}=30+6\sqrt{5}\)，故选 B.
\end{solution}

\begin{problem}
某果树前 \(n\) 年的总产量 \(S_n\) 与 \(n\) 之间的关系如图所示. 从目前记录的结果看，前 \(m\) 年的年平均产量最高，\(m\) 的值为（\quad）

\texfigure[max width=0.36\linewidth]{img_repaint/2012/beijing_liberal/q8_fig1.tex}

\choices
    {5}
    {7}
    {9}
    {11}
\end{problem}

\begin{answer}
C
\end{answer}

\begin{solution}
图中横坐标为 \(n\)，纵坐标为前 \(n\) 年总产量 \(S_n\). 因而前 \(n\) 年的年平均产量为 \(\frac{S_n}{n}\)，这正是原点 \(O\) 与图中点 \(P_n(n,S_n)\) 所连直线的斜率.

题目给出的候选值为 \(5\)，\(7\)，\(9\)，\(11\). 分别比较 \(OP_5\)、\(OP_7\)、\(OP_9\)、\(OP_{11}\) 的斜率，由图可见 \(OP_9\) 最陡，即 \(\frac{S_9}{9}\) 最大. 所以前 \(9\) 年的年平均产量最高，故 \(m=9\)，选 C.
\end{solution}

\section{填空题}

\begin{problem}
直线 \(y = x\) 被圆 \(x^2 + (y - 2)^2 = 4\) 截得的弦长为\fillinblank{}.
\end{problem}

\begin{answer}
\(2\sqrt{2}\)
\end{answer}

\begin{solution}
圆心为 \(O(0,2)\)，半径为 \(2\). 圆心到直线 \(y=x\)，即 \(x-y=0\) 的距离为
\(d=\dfrac{|0-2|}{\sqrt{1^2+(-1)^2}}=\sqrt{2}\). 设弦的中点为 \(H\)，则 \(OH\perp\) 弦，半弦长为
\(\sqrt{2^2-d^2}=\sqrt{4-2}=\sqrt{2}\)，所以弦长为 \(2\sqrt{2}\).
\end{solution}

\begin{problem}
已知 \(\{a_n\}\) 为等差数列，\(S_n\) 为其前 \(n\) 项和，若 \(a_1 = \frac{1}{2}\)，\(S_2 = a_3\)，则 \(a_2=\)\fillinblank{}；\(S_n=\)\fillinblank{}.
\end{problem}

\begin{answer}
\(1\)；\(\frac{n(n+1)}{4}\)
\end{answer}

\begin{solution}
设等差数列的公差为 \(d\)，则 \(a_2=\frac{1}{2}+d\)，\(a_3=\frac{1}{2}+2d\). 由 \(S_2=a_3\) 得
\(\frac{1}{2}+\left(\frac{1}{2}+d\right)=\frac{1}{2}+2d\)，从而 \(d=\frac{1}{2}\)，故 \(a_2=1\). 又
\(S_n=\frac{n}{2}\left[2a_1+(n-1)d\right]=\frac{n}{2}\left[1+\frac{n-1}{2}\right]=\frac{n(n+1)}{4}\).
\end{solution}

\begin{problem}
在 \(\triangle ABC\) 中，若 \(a = 3\)，\(b = \sqrt{3}\)，\(\angle A = \frac{\pi}{3}\)，则 \(\angle C\) 的大小为\fillinblank{}.
\end{problem}

\begin{answer}
\(\frac{\pi}{2}\)
\end{answer}

\begin{solution}
由正弦定理，\(\dfrac{b}{\sin B}=\dfrac{a}{\sin A}\)，所以
\(\sin B=\dfrac{b\sin A}{a}=\dfrac{\sqrt{3}\cdot\frac{\sqrt{3}}{2}}{3}=\frac{1}{2}\).
由三角形内角和 \(A+B<\pi\)，得 \(B<\frac{2\pi}{3}\)，所以在满足 \(\sin B=\frac{1}{2}\) 的两个可能值中只能取 \(B=\frac{\pi}{6}\)，不能取 \(\frac{5\pi}{6}\). 因此
\(C=\pi-A-B=\pi-\frac{\pi}{3}-\frac{\pi}{6}=\frac{\pi}{2}\).
\end{solution}

\begin{problem}
已知函数 \(f(x) = \lg x\). 若 \(f(ab) = 1\)，则 \(f(a^2) + f(b^2)=\)\fillinblank{}.
\end{problem}

\begin{answer}
\(2\)
\end{answer}

\begin{solution}
由 \(f(ab)=1\) 得 \(\lg(ab)=1\). 由于相关对数有意义，\(a^2\)，\(b^2\) 为正，且
\(f(a^2)+f(b^2)=\lg(a^2)+\lg(b^2)=\lg\left((ab)^2\right)=2\lg(ab)=2\).
\end{solution}

\begin{problem}
已知正方形 \(ABCD\) 的边长为 \(1\)，点 \(E\) 是 \(AB\) 边上的动点，则 \(\overrightarrow{DE} \cdot \overrightarrow{CB}\) 的值为\fillinblank{}；\(\overrightarrow{DE} \cdot \overrightarrow{DC}\) 的最大值为\fillinblank{}.
\end{problem}

\begin{answer}
\(1\)；\(1\)
\end{answer}

\begin{solution}
建立直角坐标系，取 \(A(0,0)\)，\(B(1,0)\)，\(C(1,1)\)，\(D(0,1)\)，设 \(E(t,0)\)，其中 \(0\leqslant t\leqslant1\). 则
\(\overrightarrow{DE}=(t,-1)\)，\(\overrightarrow{CB}=(0,-1)\)，\(\overrightarrow{DC}=(1,0)\).
因此 \(\overrightarrow{DE}\cdot\overrightarrow{CB}=1\)，且 \(\overrightarrow{DE}\cdot\overrightarrow{DC}=t\leqslant1\)，当 \(E=B\)（即 \(t=1\)）时取到最大值 \(1\).
\end{solution}

\begin{problem}
已知 \(f(x) = m(x - 2m)(x + m + 3)\)，\(g(x) = 2^{x} - 2\). 若 \(\forall x \in \R\)，\(f(x) < 0\) 或 \(g(x) < 0\)，则 \(m\) 的取值范围是\fillinblank{}.
\end{problem}

\begin{answer}
\((-4,0)\)
\end{answer}

\begin{solution}
因为 \(g(x)=2^x-2\)，所以 \(g(x)<0\) 当且仅当 \(x<1\). 因而题设条件等价于对一切 \(x\geqslant1\) 都有 \(f(x)<0\).

若 \(m=0\)，则 \(f(x)=0\)，不满足严格不等式；若 \(m>0\)，则 \(f(x)\to+\infty\)（\(x\to+\infty\)），也不可能在 \([1,+\infty)\) 上恒为负. 故必须有 \(m<0\).

此时 \(f(x)\) 是开口向下的二次函数，两个零点为 \(2m\) 和 \(-m-3\). 要使 \(x\geqslant1\) 时严格有 \(f(x)<0\)，两个零点都必须小于 \(1\). 条件 \(2m<1\) 在 \(m<0\) 时自动成立，而 \(-m-3<1\) 等价于 \(m>-4\). 所以 \(m\in(-4,0)\).
\end{solution}

\section{解答题}

\begin{problem}
已知函数 \( f(x) = \dfrac{(\sin x - \cos x)\sin 2x}{\sin x} \).
\begin{enumerate}
\item 求 \( f(x) \) 的定义域及最小正周期；
\item 求 \( f(x) \) 的单调递减区间.
\end{enumerate}
\begin{solution}
\begin{enumerate}
\item 由分母不能为零，得 \(\sin x\ne0\)，所以定义域为
\(\{x\mid x\ne k\pi,\ k\in\Z\}\). 在此定义域内，利用 \(\sin2x=2\sin x\cos x\)，得
\[
f(x)=2\cos x(\sin x-\cos x)=\sin2x-\cos2x-1
=\sqrt{2}\sin\left(2x-\frac{\pi}{4}\right)-1
\]
化简式的周期为 \(\pi\)，所以 \(\pi\) 是原函数的一个正周期. 另一方面，原函数的定义域排除了恰为 \(\{k\pi\mid k\in\Z\}\) 的点；若 \(T\) 是原函数的周期，则平移 \(T\) 必须把这组不可取点保持不变，从而 \(T\) 必须是 \(\pi\) 的整数倍. 因此原函数的最小正周期为 \(\pi\).

\item 由上式可得 \(f'(x)=2\sqrt{2}\cos\left(2x-\frac{\pi}{4}\right)\). 函数递减时
\[
\cos\left(2x-\frac{\pi}{4}\right)<0
\]
即存在 \(k\in\Z\)，使
\[
2k\pi+\frac{\pi}{2}<2x-\frac{\pi}{4}<2k\pi+\frac{3\pi}{2}
\]
解得 \(k\pi+\frac{3\pi}{8}<x<k\pi+\frac{7\pi}{8}\). 端点处导数为零，但函数在端点及其间的每一点均有定义，故可将递减区间写成相应闭区间；这些区间也不包含被排除的 \(k\pi\). 所以单调递减区间为
\(\left[k\pi+\frac{3\pi}{8},k\pi+\frac{7\pi}{8}\right]\)（\(k\in\Z\)）.
\end{enumerate}
\end{solution}
\end{problem}

\begin{problem}
如图 1，在 \(\mathrm{Rt}\triangle ABC\) 中，\(\angle C = 90^\circ\)，\(D\)，\(E\) 分别是 \(AC\)，\(AB\) 的中点，点 \(F\) 为线段 \(CD\) 上的一点，将 \(\triangle ADE\) 沿 \(DE\) 折起到 \(\triangle A_1DE\) 的位置，使 \(A_1F \perp CD\)，如图 2.

\begin{enumerate}
\item 求证：\(DE\parallel\text{平面 }A_1CB\)；
\item 求证：\(A_1F\perp BE\)；
\item 线段 \(A_1B\) 上是否存在点 \(Q\)，使 \(A_1C\perp\text{平面 }DEQ\)？说明理由.
\end{enumerate}

\texfigure[max width=0.50\linewidth]{img_repaint/2012/beijing_liberal/q16_fig1.tex}
\begin{solution}
\begin{enumerate}
\item 因为 \(D\)，\(E\) 分别是 \(AC\)，\(AB\) 的中点，所以在折叠前 \(DE\parallel BC\)，折叠后 \(DE\) 和 \(BC\) 的位置均不变，仍有 \(DE\parallel BC\). 又 \(BC\subset\text{平面 }A_1CB\). 只需证明 \(DE\) 不在该平面内.

若 \(DE\subset\text{平面 }A_1CB\)，则平面 \(A_1CB\) 与折叠前的平面 \(ABC\) 都含有两条互相平行且不重合的直线 \(DE\) 和 \(BC\)，从而两平面重合，故 \(A_1\in\text{平面 }ABC\). 由于 \(DE\perp AC\)，且 \(D\) 是 \(AC\) 的中点，点 \(A\) 绕折痕 \(DE\) 后仍落在平面 \(ABC\) 内时只能处于原位置 \(A\) 或关于 \(DE\) 的对称位置 \(C\). 这两种情形下 \(A_1F\) 均与 \(CD\) 共线或退化，不可能满足 \(A_1F\perp CD\). 因此 \(DE\not\subset\text{平面 }A_1CB\)，所以 \(DE\parallel\text{平面 }A_1CB\).

\item 由 \(AC\perp BC\) 且 \(DE\parallel BC\)，得 \(DE\perp AC\)，从而 \(DE\perp CD\). 折叠保持折痕 \(DE\) 与边 \(AD\) 的夹角，所以 \(DE\perp A_1D\). 直线 \(A_1D\) 与 \(DC\) 在 \(D\) 处相交，且都在平面 \(A_1DC\) 内，故 \(DE\perp\text{平面 }A_1DC\). 因为 \(F\in CD\)，所以 \(A_1F\subset\text{平面 }A_1DC\)，从而 \(DE\perp A_1F\).

已知 \(A_1F\perp CD\). 过 \(F\) 作直线 \(\ell\parallel DE\). 直线 \(\ell\) 和 \(CD\) 是平面 \(BCDE\) 内过 \(F\) 的两条相交直线，且 \(A_1F\perp\ell\)、\(A_1F\perp CD\)，所以 \(A_1F\perp\text{平面 }BCDE\). 由于 \(BE\subset\text{平面 }BCDE\)，故 \(A_1F\perp BE\).

\item 取 \(P\)，\(Q\) 分别为 \(A_1C\)，\(A_1B\) 的中点. 在三角形 \(A_1CB\) 中，\(PQ\parallel CB\)，而 \(DE\parallel CB\)，所以 \(PQ\parallel DE\). 平面 \(DEQ\) 含有过 \(Q\) 且平行于 \(DE\) 的直线 \(PQ\)，故 \(P\in\text{平面 }DEQ\)，从而平面 \(DEQ=\text{平面 }DEP\).

折叠保持长度，且 \(D\) 是 \(AC\) 的中点，所以 \(A_1D=AD=DC\). 因而三角形 \(A_1DC\) 是以 \(A_1C\) 为底边的等腰三角形；\(P\) 是底边 \(A_1C\) 的中点，所以 \(DP\perp A_1C\). 由上一步知 \(DE\perp\text{平面 }A_1DC\)，故 \(DE\perp A_1C\). 又 \(PQ\parallel DE\)，所以 \(PQ\perp A_1C\). 直线 \(DP\) 与 \(PQ\) 在 \(P\) 处相交，且均在平面 \(DEP\) 内，因此 \(A_1C\perp\text{平面 }DEP=\text{平面 }DEQ\). 所以存在这样的点 \(Q\)，可取 \(Q\) 为 \(A_1B\) 的中点.
\end{enumerate}
\end{solution}
\end{problem}

\begin{problem}
近年来，某市为促进生活垃圾的分类处理，将生活垃圾分为厨余垃圾，可回收物和其他垃圾三类，并分别设置了相应的垃圾箱. 为调查居民生活垃圾分类投放情况，现随机抽取了该市三类垃圾箱中总计 \(1000\text{ 吨}\) 生活垃圾，数据统计如下（单位：吨）：

\begin{center}
\begin{tabular}{c|ccc}
 & “厨余垃圾”箱 & “可回收物”箱 & “其他垃圾”箱 \\
\hline
厨余垃圾 & 400 & 100 & 100 \\
可回收物 & 30 & 240 & 30 \\
其他垃圾 & 20 & 20 & 60
\end{tabular}
\end{center}

\begin{enumerate}
\item 试估计厨余垃圾投放正确的概率；
\item 试估计生活垃圾投放错误的概率；
\item 假设厨余垃圾在“厨余垃圾”箱，“可回收物”箱，“其他垃圾”箱的投放量分别为 \(a\)，\(b\)，\(c\)，其中 \(a > 0\)，\(a + b + c = 600\). 当数据 \(a\)，\(b\)，\(c\) 的方差 \(s^2\) 最大时，写出 \(a\)，\(b\)，\(c\) 的值（结论不要求证明），并求此时 \(s^2\) 的值.
\end{enumerate}

（注：\(s^2 = \frac{1}{n}\left[(x_1 - \overline{x})^2 + (x_2 - \overline{x})^2 + \cdots + (x_n - \overline{x})^2\right]\)，其中 \(\overline{x}\) 为 \(x_1\)，\(x_2\)，\(\cdots\)，\(x_n\) 的平均数）

\begin{solution}
\begin{enumerate}
\item 三类厨余垃圾共 \(400+100+100=600\) 吨，其中正确投放到“厨余垃圾”箱的有 \(400\) 吨，所以估计正确投放的概率为 \(\dfrac{400}{600}=\dfrac{2}{3}\).

\item 投放错误的垃圾量是所有非对角线数据之和，即
\(100+100+30+30+20+20=300\) 吨. 因而估计生活垃圾投放错误的概率为 \(\dfrac{300}{1000}=\dfrac{3}{10}\).

\item 由于 \(a\)，\(b\)，\(c\) 是投放量，\(b\)，\(c\geqslant0\). 它们的平均数为 \(\bar{x}=\dfrac{a+b+c}{3}=200\)，所以
\[
s^2=\frac{1}{3}\left[(a-200)^2+(b-200)^2+(c-200)^2\right]
=\frac{a^2+b^2+c^2-120000}{3}
\]
在 \(a\)，\(b\)，\(c\geqslant0\) 且 \(a+b+c=600\) 时，
\(a^2+b^2+c^2\leqslant(a+b+c)^2=360000\)，等号当且仅当其中两个数为零. 结合 \(a>0\)，方差最大时 \(a=600\)，\(b=c=0\)，并且
\(s^2=\dfrac{360000-120000}{3}=80000\).
\end{enumerate}
\end{solution}
\end{problem}

\begin{problem}
已知函数 \(f(x) = ax^2 + 1\) （\(a > 0\)），\(g(x) = x^3 + bx\).
\begin{enumerate}
\item 若曲线 \(y = f(x)\) 与曲线 \(y = g(x)\) 在它们的交点 \((1, c)\) 处具有公共切线，求 \(a\)，\(b\) 的值；
\item 当 \(a = 3\)，\(b = -9\) 时，若函数 \(f(x) + g(x)\) 在区间 \([k, 2]\) 上的最大值为 28，求 \(k\) 的取值范围.
\end{enumerate}
\begin{solution}
\begin{enumerate}
\item 因为 \((1,c)\) 同时在两条曲线上，所以 \(a+1=c=1+b\)，从而 \(a=b\). 两条曲线在该点具有公共切线，说明导数值相等，即
\(f'(1)=2a\)，\(g'(1)=3+b\)，所以 \(2a=3+b\). 联立 \(a=b\) 得 \(a=b=3\).

\item 令 \(h(x)=f(x)+g(x)=x^3+3x^2-9x+1\)，则
\(h'(x)=3x^2+6x-9=3(x+3)(x-1)\). 因此 \(h\) 在 \((-\infty,-3)\) 上递增，在 \((-3,1)\) 上递减，在 \((1,+\infty)\) 上递增. 又
\(h(-3)=28\)，\(h(1)=-4\)，\(h(2)=3\).
区间 \([k,2]\) 要有意义，需 \(k\leqslant2\). 当 \(k\leqslant-3\) 时，区间 \([k,2]\) 包含 \(x=-3\)，且由上述单调性，\(h(x)\) 在该区间上的最大值为 \(h(-3)=28\). 当 \(-3<k\leqslant2\) 时，区间不含 \(-3\)，在 \([k,1]\) 上函数递减、在 \([1,2]\) 上函数递增，且 \(h(k)<28\)、\(h(2)=3<28\)，故最大值严格小于 \(28\). 所以 \(k\in(-\infty,-3]\).
\end{enumerate}
\end{solution}
\end{problem}

\begin{problem}
已知椭圆 \(C: \frac{x^2}{a^2} + \frac{y^2}{b^2} = 1\) （\(a > b > 0\)）的一个顶点为 \(A(2, 0)\)，离心率为 \(\frac{\sqrt{2}}{2}\)，直线 \(y = k(x - 1)\) 与椭圆 \(C\) 交于不同的两点 \(M\)，\(N\).

\begin{enumerate}
\item 求椭圆 \(C\) 的方程；
\item 当 \(\triangle AMN\) 的面积为 \(\frac{\sqrt{10}}{3}\) 时，求 \(k\) 的值.
\end{enumerate}
\begin{solution}
\begin{enumerate}
\item 由于 \(A(2,0)\) 是长轴顶点，所以 \(a=2\). 设焦半距为 \(c\)，由离心率
\(e=\frac{c}{a}=\frac{\sqrt{2}}{2}\) 得 \(c=\sqrt{2}\). 再由 \(b^2=a^2-c^2\)，得 \(b^2=4-2=2\). 所以椭圆方程为
\(\frac{x^2}{4}+\frac{y^2}{2}=1\).

\item 令 \(M(x_1,y_1)\)，\(N(x_2,y_2)\). 将 \(y=k(x-1)\) 代入椭圆方程，得到
\[
(1+2k^2)x^2-4k^2x+2k^2-4=0
\]
因此
\[
|x_1-x_2|=\frac{\sqrt{(-4k^2)^2-4(1+2k^2)(2k^2-4)}}{1+2k^2}
=\frac{2\sqrt{6k^2+4}}{1+2k^2}
\]
线段 \(MN\) 在直线 \(y=k(x-1)\) 上，所以
\(|MN|=|x_1-x_2|\sqrt{1+k^2}\). 点 \(A(2,0)\) 到该直线的距离为
\(\dfrac{|2k-k|}{\sqrt{1+k^2}}=\dfrac{|k|}{\sqrt{1+k^2}}\). 从而
\[
[\triangle AMN]=\frac{1}{2}|MN|\cdot\frac{|k|}{\sqrt{1+k^2}}
=\frac{|k|\sqrt{6k^2+4}}{1+2k^2}
\]
令 \(u=k^2\)，由题设面积得
\[
\frac{u(6u+4)}{(1+2u)^2}=\frac{10}{9}
\]
即 \(7u^2-2u-5=0\). 因而 \(u=1\) 或 \(u=-\frac{5}{7}\). 因为 \(u=k^2\geqslant0\)，所以 \(u=1\)，故 \(k=\pm1\).
\end{enumerate}
\end{solution}
\end{problem}

\begin{problem}
设 \(A\) 是如下形式的 2 行 3 列的数表，

\examdisplayarray{}{ccc}{
a & b & c \\
d & e & f
}{}

满足性质 \(P\)：\(a\)，\(b\)，\(c\)，\(d\)，\(e\)，\(f \in [-1, 1]\)，且 \(a + b + c + d + e + f = 0\).
记 \(r_i(A)\) 为 \(A\) 的第 \(i\) 行各数之和 （\(i = 1\)，\(2\)），\(c_j(A)\) 为 \(A\) 的第 \(j\) 列各数之和 （\(j = 1\)，\(2\)，\(3\)）；记 \(k(A)\) 为 \(|r_1(A)|\)，\(|r_2(A)|\)，\(|c_1(A)|\)，\(|c_2(A)|\)，\(|c_3(A)|\) 中的最小值.

\begin{enumerate}
\item 对如下数表 \(A\)，求 \(k(A)\) 的值；

\examdisplayarray{}{ccc}{
1 & 1 & -0.8 \\
0.1 & -0.3 & -1
}{}

\item 设数表 \(A\) 形如

\examdisplayarray{}{ccc}{
1 & 1 & -1 - 2d \\
d & d & -1
}{}

其中 \(-1 \leqslant d \leqslant 0\). 求 \(k(A)\) 的最大值；
\item 对所有满足性质 \(P\) 的 2 行 3 列的数表 \(A\)，求 \(k(A)\) 的最大值.
\end{enumerate}
\begin{solution}
\begin{enumerate}
\item 该数表的两行和分别为 \(1.2\) 和 \(-1.2\)，三列和分别为 \(1.1\)，\(0.7\) 和 \(-1.8\). 所以
\(k(A)=\min\{1.2,1.2,1.1,0.7,1.8\}=0.7\).

\item 该数表的行和、列和分别为
\(r_1=1-2d\)，\(r_2=-1+2d\)，\(c_1=c_2=1+d\)，\(c_3=-2-2d\). 因为 \(-1\leqslant d\leqslant0\)，有
\(|r_1|=|r_2|=1-2d\geqslant1+d\geqslant0\)，且 \(|c_3|=2(1+d)\geqslant1+d\). 因此 \(k(A)=1+d\)，其最大值在 \(d=0\) 时取得，最大值为 \(1\).

\item 交换两行、交换三列以及同时将全部元素变号，都不改变 \(k(A)\). 三个列和的总和为零，所以至少有两个列和同为非负数或同为非正数；通过交换列并在必要时将全部元素变号，可设所选两列为 \(c_1\geqslant0\)、\(c_2\geqslant0\). 再交换两行，可设 \(r_1\geqslant0\). 因而
\(k(A)\leqslant r_1\)，\(k(A)\leqslant c_1\)，\(k(A)\leqslant c_2\)，从而
\[
3k(A)\leqslant r_1+c_1+c_2
=(a+b+c)+(a+d)+(b+e)
=a+b-f\leqslant3
\]
所以 \(k(A)\leqslant1\). 第（2）问取 \(d=0\) 时，数表仍满足性质 \(P\)，且 \(k(A)=1\)，故所有满足性质 \(P\) 的数表中 \(k(A)\) 的最大值为 \(1\).
\end{enumerate}
\end{solution}
\end{problem}
